\chapter{Conclusion}
\label{chap:conclusion}

This thesis investigated whether a learned embedding-based heuristic can make $A^*$ a practical approach for binary causal question answering over large causal knowledge graphs. Questions of the form ``Does $A$ cause $B$?'' were modeled as directed reachability problems, with fine-tuned embedding distances guiding graph traversal. The final configuration uses Granite Embedding R2 with a 32-dimensional Matryoshka prefix, ReLU activation, Euclidean distance, and a search budget of $\tau=27$. The results show that embedding-guided $A^*$ substantially reduces traversal cost while maintaining competitive predictive performance. Its efficiency advantage is particularly pronounced on large or dense graphs such as CauseNet Full and CEG Filtered. On these graphs, $A^*$ is at least $9.8\times$ faster than capped BFS and $32.1\times$ faster than uncapped BFS, while visiting hundreds to thousands of times fewer nodes. It also outperforms the evaluated RL baseline in F\textsubscript{1} and runtime across all six final test settings. Overall, the findings demonstrate that learned embedding guidance provides an efficient alternative to exhaustive and policy-based traversal for binary causal graph search. Future work could address the limitations identified in this thesis. Adaptive, query-dependent search budgets could reduce recall loss caused by bounded exploration. Incorporating edge confidence, provenance, and path-quality estimation could improve robustness to noisy or ambiguous graph relations. Further work could also investigate the contextual validity of composed causal paths and extend the approach beyond binary reachability to tasks involving causal relation quality, effect strength, alternative causes, and interventional or counterfactual reasoning.